\chapter{Conclusiones}
\label{sec.ej3:conclusiones}

El desarrollo experimental llevado a cabo sobre el Kit de Comunicaciones N° 5 permitió analizar y contrastar en forma
directa las características temporales y espectrales de las modulaciones DSB-SC y AM convencional utilizando el
integrado modulador balanceado MC1496. A partir de los resultados obtenidos y las mediciones efectuadas con el
osciloscopio digital, se destacan las siguientes conclusiones técnicas:

\begin{itemize}
    \item \textbf{Eficiencia de potencia y distribución espectral:} En la modulación AM estándar, la potencia
          transmitida está dominada por la portadora de radiofrecuencia. Incluso bajo la condición óptima de modulación
          completa ($m = 1{,}0$), dos tercios ($66{,}7\%$) de la potencia total emitida corresponden a la portadora
          pura, la cual no transporta información en banda base; este desperdicio energético asciende al $88{,}9\%$
          cuando el índice de modulación se reduce a $m = 0{,}5$. En cambio, en la modulación DSB-SC la portadora
          resulta suprimida en el origen, concentrando el $100\%$ de la potencia de emisión en las bandas laterales de
          información (USB y LSB), lo que representa una ganancia sustancial de eficiencia para enlaces de potencia
          restringida.
          
    \item \textbf{Compromiso entre eficiencia y complejidad de recepción:} Si bien DSB-SC aventaja a la AM estándar
          en rendimiento energético, exige una mayor complejidad circuital en la etapa receptora. Mientras que la AM
          permite una demodulación asincrónica sumamente sencilla mediante detección de envolvente por diodo, la señal
          DSB-SC requiere forzosamente detección sincrónica o coherente. Dicha técnica exige regenerar localmente una
          portadora en estricta sincronía de frecuencia y fase respecto a la del modulador, dado que cualquier error
          angular $\phi$ atenúa la señal recuperada según la ley $\cos(\phi)$, anulándola completamente ante una
          condición de cuadratura ($\phi = 90^\circ$).

    \item \textbf{Dinámica temporal e inversión de fase:} Las capturas en el osciloscopio evidenciaron claramente la
          diferencia fundamental en el dominio temporal: mientras que en AM la envolvente replica directamente la señal
          modulante senoidal sin cambio de fase (siempre que $m \le 1$), en DSB-SC la envolvente se anula periódicamente
          en cada cruce por cero de la modulante, registrándose una inversión instantánea de fase de $180^\circ$ en la
          portadora interna.

    \item \textbf{Régimen de sobremodulación en AM:} Al forzar una condición de sobremodulación ($m \approx 1{,}5$)
          inyectando una modulante que superó el nivel continuo de portadora ($E_m > E_c$), se constató la aparición de
          cruces espurios por el eje de referencia y recortes en la envolvente. En el análisis espectral por FFT, esto se
          tradujo en una severa distorsión no lineal con aparición de productos armónicos laterales de orden superior
          en $f_c \pm 2 f_m$ y $f_c \pm 3 f_m$. Este ensanchamiento espectral corrobora el riesgo de interferencia
          hacia canales de RF contiguos (splatter), justificando los límites estrictos impuestos a los transmisores
          comerciales.

    \item \textbf{Ancho de banda ocupado:} En ambos esquemas de modulación lineal, las mediciones con los cursores de
          la FFT demostraron rigurosamente que el ancho de banda ocupado en radiofrecuencia responde a la relación
          teórica $B = 2 f_m$. Variando la modulante desde $1\kHz$ hasta $20\kHz$, el canal pasó de ocupar $2\kHz$
          a $40\kHz$, manteniendo la estricta simetría respecto a la frecuencia central de portadora de $100\kHz$.

    \item \textbf{Flexibilidad operativa del integrado MC1496:} La celda multiplicadora de cuatro cuadrantes
          demostró una notable versatilidad. Mediante el preset de polarización estática de los emisores diferenciales,
          fue posible controlar la reinserción de portadora en un rango dinámico de $5{,}90\dB$ para AM estándar, así
          como balancear con alta precisión las ramas del circuito sin señal modulante, alcanzando un factor de
          rechazo de portadora de $13{,}46\dB$ para operar en modo DSB-SC.
\end{itemize}

En síntesis, los ensayos de laboratorio permitieron validar de manera integral los conceptos teóricos de modulación y
demodulación lineal de amplitud, adquirir destreza en el manejo del instrumental digital de medición y comprender las
ventajas, desventajas y criterios de diseño asociados a cada esquema en los sistemas de comunicaciones analógicas.
